% ADASS 2026 proceedings, paper P301
\documentclass[11pt,twoside]{article}

% Do NOT use ANY packages other than asp2023.
\usepackage{asp2023}

\aspSuppressVolSlug
\resetcounters

\bibliographystyle{asp2021}

\markboth{Salgado}{Use of AI for SRCNet Development}

\begin{document}

\title{Use of AI for SRCNet Development}

\author{Jes\'us~Salgado}
\affil{SKA Observatory, Jodrell Bank, Lower Withington, Macclesfield, UK; \email{jesus.salgado@skao.int}}

\paperauthor{Jes\'us~Salgado}{jesus.salgado@skao.int}{}{SKA Observatory}{}{Macclesfield}{Cheshire}{SK11 9FT}{United Kingdom}

%\aindex{Salgado,~J.}

\begin{abstract}
The SKA Regional Centre Network (SRCNet) is preparing for SKA-Low Science
Verification, with tight milestones and many interdependent services that must
mature quickly and then be maintained for years. We report on the use of AI
coding assistants and large language models (LLMs) across the SRCNet
development teams. Our approach relies on structured prompting, in which each
request encodes the target architecture, module boundaries, documentation
standards and tests from the outset, followed by an early handoff of the
prototype to the team that will own it. We describe the governance framework
at the SKA Observatory, four ways in which LLMs are used in SRCNet, and the
practices adopted to keep human review meaningful when code can be generated
faster than it can be read. Early evidence from our GitLab and SonarQube
instances suggests that, managed this way, AI-assisted services are delivered
faster without a loss of measured quality.
\end{abstract}

%\ssindex{Artificial intelligence}
%\ssindex{observatories!ground-based!SKA}
%\ssindex{algorithm!machine learning}
%\ssindex{software}
%\ssindex{quality control}
%\ooindex{astroquery, ascl:1708.004}

\section{Introduction}

The Square Kilometre Array \citep{2009IEEEP..97.1482D} will deliver around
600\,PB of science data products per year to a federated network of regional
centres, the SRCNet \citep{2024ASPC..535..399S, 2025ASPC..541...64S}. The
network is built by about 70 people in 16 countries, organised in several
development teams, and is now on the road to Science Verification
\citep{2025ASPC..541..385C}: registration, ingestion, metadata, data access
over a Rucio data lake \citep{2019CSBS....3...11B}, federated execution and
accounting all need to mature at the same time.

LLMs trained on code \citep{2021arXiv210703374C} have become everyday tools.
Controlled studies report large speed-ups on bounded tasks
\citep{2023arXiv230206590P}, but also slowdowns for experienced developers
working on mature code they know well \citep{2025arXiv250709089B}, and
industry surveys describe AI as an amplifier of a team's existing practices,
good or bad \citep{DORA2025}. For SRCNet, the question was therefore not
whether to use AI, but how to use it without producing code that nobody can
review or maintain.

\section{Governance}

The SKAO AI policy \citep{SKAO-GOV-0000166} rates every AI use on a
five-level risk scale. Code generation and chatbots are moderate risk:
generated code follows the SKA quality-assurance process, with review and
testing, and must be verified and understood before use. Only Unrestricted
information (source code and technical documentation) may be shared with
hosted tools. The SKAO adoption plan \citep{SKAO-GOV-0000194} adds an opt-in,
three-tier user model and a community of practice. In SKAO developer surveys,
weekly AI use rose from 62\% (2025) to 86\% (2026), and reported productivity
gains from 57\% to 78\%.

\section{Four Uses of LLMs in SRCNet}

\begin{description}
\item[Code generation (in use).] Services, clients, notebooks and CI are
prototyped with an AI assistant and handed to the teams, e.g.\ the
\texttt{astroquery.srcnet} client \citep{2019AJ....157...98G}: four
repositories, five notebooks and a three-stage CI from one project brief.
\item[Reports and artefacts (in use).] Reviews, interface documents and
presentations are drafted only from source documents and checked by a person;
grounded drafting also surfaces inconsistencies between documents and code.
\item[Hiding logic behind language (prototype).] A chat service turns
plain-language questions into ADQL \citep{2023ivoa.spec.1215M} run over TAP
\citep{2019ivoa.spec.0927D} with the user's own token; the model is swappable
and can be self-hosted.
\item[Agents for operations (under study).] Agents that read Rucio and
transfer logs and propose actions, behind a sanitisation step, with read-only
tools and a human approving every action.
\end{description}

\section{Structured Prompting and Handoff}

The work follows five steps: \emph{frame} (an architect writes a shared
project brief and a structured prompt), \emph{prototype} (the assistant stands
up skeletons, explores API contracts and validates design choices),
\emph{hand off} (the prototype, its design decisions and its documentation
become a feature in a team's backlog), \emph{harden} (integration, SKAO CI/CD,
at least 80\% coverage, SonarQube and review) and \emph{maintain}
(by people, with or without AI). Every prompt encodes four
things: the target architecture and design patterns (for example a singleton
entry point and factory methods), module boundaries and interfaces,
documentation for the next human maintainer, and tests written with the code
rather than after it.

\section{The Review Bottleneck and Early Evidence}

An assistant can write thousands of lines in one session, but a reviewer cannot
judge them in one sitting, and review effectiveness is known to depend on
understanding the change \citep{Bacchelli2013}. Oversized merge requests risk
becoming a bottleneck or a rubber stamp. We mitigate this with prompt-scoped
increments (one concern per merge request), review checklists aligned to the
architecture, tests generated with the code, and documentation for the next
maintainer.

Table~\ref{tab:mr} compares merge requests labelled as AI-assisted in the
public SRCNet GitLab group with all other merge requests in 2026.
AI-assisted work grew sharply in 2026 (Figure~\ref{fig:f1}a). Those merge
requests are about twice as large and are merged much faster, with the same
median number of human comments and far more automated AI review. CI pass
rates and the share of added lines that are tests are unchanged. In SonarQube,
five core services built AI-first (the software artefact importer, the
registration orchestrator, two accounting services and the computing broker)
compare well with 27 other SRCNet projects (Figure~\ref{fig:f1}b). These are
correlations, not causes: the team grew and a quality campaign started over the
same period, the AI-first services are young, and unlabelled AI-assisted work
is not counted.

\begin{table}[!ht]
\caption{AI-labelled merge requests in 2026 against all other merge requests in the SRCNet GitLab group, to 23 October 2026.}
\label{tab:mr}
\smallskip
\begin{center}
{\small
\begin{tabular}{lcc}
\tableline
\noalign{\smallskip}
Metric & AI-labelled & Other \\
\noalign{\smallskip}
\tableline
\noalign{\smallskip}
Median lines changed & 166 & 75 \\
Merge requests over 1000 lines & 25\% & 7.5\% \\
Median time to merge & 2.5\,h & 22\,h \\
Median human comments & 2 & 2 \\
Reviewed by an AI review bot & 84\% & 41\% \\
CI pipeline passing at merge & 94\% & 93\% \\
\noalign{\smallskip}
\tableline
\end{tabular}
}
\end{center}
\end{table}

\articlefigure[width=0.92\textwidth]{P301_f1.png}{fig:f1}{(a) Merge requests labelled as AI-assisted per quarter in the SRCNet GitLab group (Q4 2026 to 23 October). (b) SonarQube medians per project for five services built AI-first (dark) and 27 other SRCNet projects (hatched).}

\section{Observations and Conclusions}

AI assistance has been most effective for greenfield prototypes, service
skeletons, documentation and reports, where a brief and a structured prompt fix
the shape of the result in advance. Human judgement remains indispensable for
module boundaries, cross-team interfaces, security and the decision to merge.
Two risks need attention: large, quickly merged changes, and AI-assisted work
concentrated in a few developers. Next, we will label every AI-assisted merge
request and track its size, review time and quality per quarter. Managed this
way, AI speeds up SRCNet development without giving up review or
maintainability: it amplifies good engineering practice rather than replacing it.

\acknowledgements The author thanks the SRCNet development teams and the SKAO Information Security and software quality teams for the practices and data reported here.

\bibliography{P301}

% if we have space left, we might add a conference photograph here. Leave commented for now.
% \bookpartphoto[width=1.0\textwidth]{foobar.eps}{FooBar Photo (Photo: Any Photographer)}

\end{document}
